\apartado{2.2}{Tipos de probabilidad}{fundamentos}
\bajada{Tres maneras de ponerle un número a lo incierto. En psicología las
tres tienen lugar propio, y no son intercambiables: una la fabrica el diseño,
otra sale de contar, y la tercera es un juicio profesional.}

\definicion{Los tres tipos}{%
  \textbf{Clásica}: $P(A)=\text{favorables}/\text{posibles}$, \emph{sólo} si
  todos los resultados son igualmente probables.\par
  \textbf{Frecuentista}: la proporción observada en muchos casos.\par
  \textbf{Subjetiva}: un grado de creencia fundado, sin conteo detrás.}

\minihistoria{Dónde vive cada una en tu trabajo}{%
  \textbf{Clásica} — asignás 40 participantes al azar a tratamiento o control.
  La probabilidad de que a alguien le toque tratamiento es exactamente
  $1/2$, y lo sabés \emph{antes} de empezar: la simetría la fabricaste vos con
  el procedimiento.\par
  \textbf{Frecuentista} — de 200 pacientes que iniciaron terapia, 46
  abandonaron antes de la sexta sesión. La tasa de abandono es $46/200$, y no
  se podía saber de antemano: hubo que observar.\par
  \textbf{Subjetiva} — un terapeuta con veinte años de práctica dice que este
  paciente tiene buen pronóstico. No hay conteo detrás, y sin embargo es
  información valiosa.}

\ejercicio{0}{¿Cuáles de estos números pueden ser una probabilidad? Marcá los
  que no y explicá por qué.
  \[ -0{,}2 \qquad 0{,}5 \qquad 1{,}4 \qquad \tfrac{3}{4}
     \qquad 0 \qquad 1 \qquad 46/200 \]}{}
\renglones[3]
\respuesta{Válidos: $0{,}5$, $3/4$, $0$, $1$ y $46/200$. No: $-0{,}2$ y $1{,}4$.}
\solucion{%
  $-0{,}2$ no puede serlo porque ninguna probabilidad es negativa, y $1{,}4$
  tampoco porque ninguna supera a 1.
  \par
  Que $0$ y $1$ sí valgan no es un detalle técnico: $0$ es el evento imposible
  y $1$ el evento seguro. Y $46/200=0{,}23$ está perfectamente dentro del
  rango.
  \par
  Estos dos límites son los dos primeros axiomas de Kolmogórov, y son la única
  forma rápida de detectar que una cuenta salió mal: si te dio $1{,}3$, no hace
  falta revisar nada más — ya sabés que hay un error.}
\enclase{Pediles que digan en voz alta qué significaría una probabilidad de
$1{,}4$. Nadie puede completar la frase, y esa imposibilidad enseña más que
el axioma escrito.}

\ejercicio{1}{Un estudio asigna al azar 40 participantes: 20 a tratamiento y
  20 a control.
  \textbf{Antes de calcular, escribí qué probabilidad creés que tiene un
  participante de quedar en tratamiento:} \underline{\hspace{2cm}}
  \par\vspace{0.2em}
  Ahora calculala. ¿Qué tipo de probabilidad usaste y por qué te alcanzó con
  el enunciado, sin observar ningún dato?}{Predecí primero}
\renglones[4]
\respuesta{$20/40 = 0{,}5$. Clásica: la simetría la fabricó el diseño.}
\solucion{%
  $P(\text{tratamiento}) = \dfrac{20}{40} = 0{,}5$.
  \par
  Es \textbf{clásica}, y ésta es la parte que importa: no hizo falta correr el
  estudio ni observar a nadie. La probabilidad se conoce de antemano porque el
  investigador \emph{construyó} la simetría al asignar al azar en partes
  iguales.
  \par
  Es de los pocos lugares en psicología donde la probabilidad clásica se
  aplica con total legitimidad — y no es casualidad: se aplica justamente
  porque alguien la diseñó para que se aplicara.}

\ejercicio{1}{De 200 pacientes que iniciaron terapia en el año, 46 abandonaron
  antes de la sexta sesión. Calculá la probabilidad de que un paciente elegido
  al azar de ese grupo haya abandonado, y la de que no.}{}
\renglones[3]
\respuesta{$46/200 = 0{,}23$ y $154/200 = 0{,}77$.}
\solucion{%
  $P(\text{abandono}) = \dfrac{46}{200} = 0{,}23$.
  \par
  Para el complemento no hace falta volver a contar:
  $P(\text{no abandono}) = 1 - 0{,}23 = 0{,}77$.
  \par
  Usar el complemento en vez de recontar ahorra trabajo y evita errores. En
  2.5 va a ser decisivo.}

\ejercicio{2}{Dos situaciones que dan el \emph{mismo número} y no son la misma
  cosa. Calculá las dos y explicá en qué se diferencian.
  \begin{enumerate}[label=(\alph*),topsep=2pt,itemsep=1pt]
    \item Probabilidad de que a un participante le toque el grupo control, en
      una asignación al azar de 50 y 50.
    \item Proporción de pacientes que completaron el tratamiento, si de 100
      lo completaron 50.
  \end{enumerate}}{Dos casos casi iguales}
\renglones[5]
\respuesta{Las dos dan $0{,}5$, pero (a) es clásica y se sabe de antemano;
(b) es frecuentista y hubo que observar.}
\solucion{%
  Las dos dan $0{,}5$. La diferencia no está en el número sino en
  \textbf{de dónde salió}:
  \par
  (a) Se sabía \textbf{antes} de empezar, porque el diseño lo garantiza. Si
  repetís el estudio, sigue siendo $0{,}5$ exactamente.
  \par
  (b) Se supo \textbf{después} de observar 100 casos. Si seguís a los próximos
  100 pacientes, casi seguro no da 50 de nuevo — puede dar 47 o 54.
  \par
  \textbf{Por qué importa}: sobre (a) podés apostar sin datos; sobre (b) tenés
  que preguntarte cuántos casos observaste. Cincuenta de 100 y 5 de 10 dan el
  mismo $0{,}5$ y no merecen la misma confianza.}
\enclase{Preguntá cuál de las dos cambiaría si repitieran el estudio. Es la
forma más rápida de que vean que el mismo número puede tener dos estatutos
distintos.}

\ejercicio{2}{Una investigadora escribe:
  \begin{quote}\itshape
  «El cuestionario tiene 28 puntajes posibles y tres niveles de gravedad
  (leve, moderado, grave). Como hay tres niveles, la probabilidad de que un
  paciente caiga en el nivel grave es $1/3$.»
  \end{quote}
  \textbf{Señalá el paso exacto} donde se equivoca y explicá por qué.}{Encontrá el error}
\renglones[5]
\respuesta{Aplica la fórmula clásica sin que haya simetría: los tres niveles no
son igualmente probables. Habría que contar.}
\solucion{%
  El error está en el paso «como hay tres niveles, cada uno vale $1/3$».
  \par
  Contar \emph{categorías} no autoriza a repartir la probabilidad en partes
  iguales entre ellas. Los tres niveles de gravedad cubren rangos distintos de
  puntaje y, sobre todo, la gente no se distribuye de manera uniforme: la
  mayoría puntúa bajo y muy pocos llegan a grave.
  \par
  Es el mismo error del $1/28$ del apartado anterior, con otro disfraz. Y es
  más peligroso acá, porque «tres niveles, un tercio cada uno» suena razonable.
  \par
  \textbf{Lo correcto}: contar cuántos pacientes del archivo caen en cada
  nivel y dividir por el total.}

\ejercicio{3}{De las 200 fichas del servicio, 43 llegan al corte de 10 o más y
  25 tienen diagnóstico confirmado por un profesional. Se elige una ficha al
  azar.
  \begin{enumerate}[label=(\alph*),topsep=2pt,itemsep=1pt]
    \item Calculá $P(\text{dar positivo})$ y $P(\text{no dar positivo})$.
    \item Calculá $P(\text{tener el diagnóstico})$. Ese número tiene nombre
      propio: ¿cuál?
    \item ¿Qué tipo de probabilidad son, y por qué no la clásica?
  \end{enumerate}}{}
\renglones[6]
\respuesta{(a) $0{,}215$ y $0{,}785$. (b) $0{,}125$, la prevalencia.
(c) Frecuentista.}
\solucion{%
  (a) $P = \dfrac{43}{200} = 0{,}215$, y su complemento $0{,}785$.
  \par
  (b) $P = \dfrac{25}{200} = 0{,}125$. Se llama \textbf{prevalencia}: qué
  proporción de la población tiene la condición. Guardalo — en 2.6 va a ser la
  mitad del cálculo que explica por qué un test bueno da tantas falsas alarmas.
  \par
  (c) Son \textbf{frecuentistas}: salen de contar lo que ocurrió en 200
  personas. No pueden ser clásicas porque nada garantiza que los puntajes sean
  igualmente probables. Nadie diseñó esa simetría, y la naturaleza no la
  regala.}

\ejercicio{3}{El servicio quiere estimar la probabilidad de que un estudiante
  con diagnóstico confirmado \emph{responda} a la primera línea de tratamiento.
  Tiene el archivo de 200 fichas con puntajes, diagnóstico y fecha de la
  entrevista. \textbf{¿Alcanza para responder?} Si sí, explicá cómo; si no,
  decí qué falta.}{¿Alcanzan los datos?}
\renglones[4]
\respuesta{No alcanza: el archivo no registra si recibieron tratamiento ni cómo
respondieron.}
\solucion{%
  \textbf{No alcanza.} El archivo llega hasta el diagnóstico, y la pregunta es
  sobre lo que pasa \emph{después}: haría falta saber quiénes iniciaron
  tratamiento y con qué resultado.
  \par
  Los 25 con diagnóstico confirmado están ahí y son tentadores, pero
  responderían otra pregunta.
  \par
  \textbf{La regla que conviene fijar}: una probabilidad frecuentista se
  calcula sobre casos donde el evento pudo observarse. Si el dato nunca se
  registró, no hay cuenta que lo produzca — hay que ir a buscarlo.}

\ejercicio{4}{Clasificá cada afirmación como \emph{clásica}, \emph{frecuentista}
  o \emph{subjetiva}, y justificá en una línea.
  \begin{enumerate}[label=(\alph*),topsep=2pt,itemsep=1pt]
    \item «A cada participante le toca el grupo experimental con probabilidad $1/2$.»
    \item «En nuestro servicio, el 23\% de los pacientes abandona antes de la sexta sesión.»
    \item «Creo que esta paciente tiene un 70\% de chances de sostener el tratamiento.»
    \item «La probabilidad de recaída al año es del 40\%.»
  \end{enumerate}
  Después contestá: ¿cuál de las cuatro cambiaría si el estudio se repitiera
  con otras personas?}{}
\renglones[6]
\respuesta{(a) clásica; (b) frecuentista; (c) subjetiva; (d) frecuentista.
Cambiarían (b) y (d).}
\solucion{%
  (a) \textbf{Clásica}: la simetría está fabricada por el diseño.\par
  (b) \textbf{Frecuentista}: sale de contar casos observados.\par
  (c) \textbf{Subjetiva}: es un juicio clínico. No hay conteo detrás, lo que no
  la vuelve inútil — es información de alguien con experiencia — pero no se
  verifica del mismo modo.\par
  (d) \textbf{Frecuentista}: sale de estudios de seguimiento.
  \par
  \textbf{Cambiarían (b) y (d)}, porque dependen de a quién se observó. La (a)
  no cambia nunca, porque no depende de las personas sino del procedimiento.
  Y la (c) cambia con el clínico, no con los pacientes — que es una tercera
  cosa, y vale la pena decirlo en voz alta.
  \par
  Las tres respetan los mismos axiomas y dan números entre 0 y 1. Lo que cambia
  es de dónde sale el número, y por lo tanto cuánto se puede confiar en él y
  qué haría falta para discutirlo.}
\enclase{La pregunta final es la que separa a quien clasificó de memoria de
quien entendió. Dejala para el cierre de la clase.}

\trampa{decir «es subjetiva» como si fuera un descalificativo}%
  {suena a «opinión sin fundamento», y en un contexto científico eso parece
   malo}%
  {una probabilidad subjetiva bien calibrada, de un profesional con
   experiencia, puede ser mejor que una frecuentista calculada sobre 10 casos.
   Lo que la distingue no es su calidad sino cómo se la discute: no se
   verifica contando, se verifica siguiendo pronósticos en el tiempo}

\cierre{%
  Los tres tipos dan un número entre 0 y 1 y respetan los mismos axiomas.
  Cambia el origen: la clásica la fabrica el diseño, la frecuentista sale de
  contar, la subjetiva es un juicio.
  \par
  Y quedó algo pendiente que no se puede resolver con un solo número: hasta
  acá cada probabilidad miraba \emph{una} variable por vez. Las preguntas que
  de verdad importan en la clínica cruzan dos — lo que dijo el test contra lo
  que era verdad, lo que dijo un evaluador contra lo que dijo el otro. Para eso
  hace falta una tabla.}
